\begin{table*}[t]
\centering
\small
\setlength{\tabcolsep}{6pt}
\textbf{(a) Task and evaluation scope}\\[3pt]
\begin{tabular}{llllc}
\toprule
Resource & Primary input & Reported scope & Main unit & Hierarchy \\
\midrule
Open Problems & Multiple & 12 living tasks & Task-specific & Task-specific \\
scMultiBench & RNA/ADT/ATAC & 64 real + 22 simulated & Cell/dataset & No \\
SPARE & scRNA-seq & 5 datasets, 8 methods & Sample & No \\
PaSCient & scRNA-seq & 12.5M cells, 2,700 patients & Patient & No \\
APT-Bench & 293 aptamers & 1 cohort, 40 patients & Cell + patient & Yes \\
\bottomrule
\end{tabular}
\vspace{6pt}

\textbf{(b) Protocol components and availability}\\[3pt]
\begin{tabular}{lcccc}
\toprule
Resource & Patient task & Two-axis scale & Assay budget & Public pipeline \\
\midrule
Open Problems & Task-specific & No & No & Yes \\
scMultiBench & No & No & Feature selection & Yes \\
SPARE & Yes & No & No & Yes \\
PaSCient & Yes & No & No & Not audited \\
APT-Bench & Secondary & Separate report & Secondary & Pending release \\
\bottomrule
\end{tabular}
\caption{Scope comparison with recent single-cell evaluation resources
\citep{luecken2025open,liu2025multitask,shitov2025spare,liu2026pascient}.
``Task-specific'' means that the living platform defines multiple contracts
rather than one shared clinical protocol. APT-Bench is narrower in cohorts but
centers a fixed aptamer modality and patient-disjoint coarse-to-fine cell
typing. Patient-level analysis is secondary; sampling studies are in a separate
report and do not expand the main contribution. Public release remains a submission requirement, not a completed
claim.}
\label{tab:benchmark_comparison}
\end{table*}
